\providecommand{\toplevelprefix}{../..}  %
\documentclass[../../book-main_es.tex]{subfiles}

\begin{document}

\chapter*{Prefacio a la versión 2.0}

% \begin{center}
%     «{\em Todos los caminos conducen a Roma}.»

% \end{center}
% \vspace{5mm}

Como libro de código abierto, la primera versión, titulada «{\em Aprendizaje de representaciones 
profundas de distribuciones de datos}», lanzada el 18 de agosto de 2025, se ha actualizado frecuentemente en los últimos seis meses. 
Al mismo tiempo , basados en nuestra experiencia usando el libro en un nuevo curso impartido 
en la Universidad de Hong Kong en el semestre de otoño de 2025 y a los comentarios que hemos 
recibido de colegas, tutores y estudiantes, hemos identificado numerosos puntos a lo largo 
del libro que merecen mayor revisión y ampliación.

Por ello hemos decidido introducir algunos cambios y actualizaciones sustanciales en el contenido 
y la organización del libro para la nueva \textit{Versión 2.0} y hemos cambiado su título a:
\begin{quote}
    \centering 
    {\em Principios y práctica del aprendizaje profundo de representaciones} \\
    {\em o Una teoría matemática de la memoria}
\end{quote}
Esta nueva versión también nos permite poner de relieve las fuertes conexiones conceptuales y 
técnicas entre los contenidos de los diferentes capítulos y secciones del libro, de manera 
que creemos que el valor pedagógico global del libro ha mejorado significativamente 
en relación a la primera versión. 

En nuestra opinión, la nueva versión logra una presentación mucho más unificada, completa y 
fluida del tema. La nueva versión también incorpora nuevos aportes teóricos recientemente desarrollados, 
así como un número creciente de aplicaciones prácticas para datos del mundo real que no se han documentado 
de manera apropiada ni explicado en forma sistemática en otros lugares. Esta Versión 2.0 servirá como 
remedio oportuno para esta situación.

\subsection*{Principales cambios en la Version 2.0}
\begin{itemize}
    \item Hemos dividido en dos capítulos el capítulo 3 de la primera versión, que ahora
        son los capítulos \ref{ch:denoising} y \ref{ch:lossy-compression}. El
        nuevo capítulo \ref{ch:denoising} se centra en el proceso de reducción 
        de ruido para el aprendizaje de distribuciones de baja dimensión.
        Hemos agregado una nueva sección en la que se caracterizan las 
        condiciones bajo las que el proceso conduce a una generalización o
        memorización. El nuevo capítulo \ref{ch:lossy-compression} se centra
        en el aprendizaje de representaciones basado en un enfoque de 
        codificación con pérdidas y en el principio de maximización de la 
        ganancia de información. Como la versión anterior solo
        ilustraba cómo aplicar el principio en el aprendizaje de 
        representaciones de imágenes en una configuración supervisada, en la nueva versión 
        se agrega un estudio de casos correspondientes a una configuración no supervisada.
        También proporciona una justificación teórica de algunos métodos populares 
        de aprendizaje no supervisado tales como el aprendizaje contrastivo 
        y DINO en particular.
    \item En el capítulo \ref{ch:deep-networks} actualizado sobre el diseño
        de representaciones por vía de optimización desenrollada, hemos agregado
        una derivación basada en principios para una versión causal de la arquitectura 
        de transformador de white-box CRATE, que resulta importante para el
        procesamiento de datos secuenciales tales como textos, como vemos
        frecuentemente en aplicaciones.
    \item En el capítulo \ref{ch:consistent-representations} actualizado 
        sobre el aprendizaje de representaciones consistentes, hemos agregado 
        una nueva sección que aumenta la claridad sobre las relaciones entre
        el aprendizaje de distribuciones y el aprendizaje de representaciones. Esto da
        una clarificación y justificación teóricas para los métodos populares
        y prácticos basados en autocodificación, como la autocodificación
        variacional o los autocodificadores de representaciones. 
    \item En el capítulo \ref{ch:conditional-inference} actualizado hemos
        añadido una sección que ofrece una explicación basada en principios
        del aprendizaje de representaciones con datos pareados y generación 
        condicional a través de la perspectiva de la infrmación mutua. Proporciona
        una justificación teórica para métodos populares y prácticos tales como
        CLIP y formas de atención cruzada.
    \item El capítulo \ref{ch:applications} ampliado sobre aplicaciones ahora
        incluye numerosas implementaciones detalladas del aprendizaje de
        distribuciones y de representaciones para de datos reales,
        incluyendo imágenes idimensionales2D, objetos naturales tridimensionales, 
        movimientos corporales, así como lenguajes naturales, en
        casi todos las configuraciones populares y prácticas, supervisadas, débilmente 
        supervisadas, no supervisadas y condicionadas.
    \item El capítulo \ref{ch:future} final, sobre líneas de 
        investigación abiertas, también se ha reescrito significativamente.
        Buscamos ofrecer una taxonomía más clara de los diferentes niveles
        de la inteligencia, de modo de poder clarificar qué hemos hecho y qué
        entendemos acerca de la inteligencia (con este libro) y qué no hemos hecho. Creemos
        que los problemas abiertos asociados con las formas superiores de
        inteligencia pueden formularse adecuadamente de manera de que puedan 
        estudiarse cualitativamente y cuantitativamente con de medios
        científicos y matemáticos, en lugar de permanecer como un asunto
        misterioso o meramente filosófico.
\end{itemize}

\subsection*{Contribuidores a la Versión 2.0}
Además de los autores, numerosos estudiantes y colegas han colaborado a este proyecto y han aportado contenido valioso a diferentes partes del libro durante la preparación de la {\em Versión 2.0}. Se presenta a continuación una lista incompleta de las personas y sus contribuciones específicas, en orden alfabético:
\begin{itemize}
    \item Tianzhe Chu: Experimentos en la Sección \ref{sec:image_classification}, así como herramientas de IA para el libro.
    \item Prof. Shenghua Gao: Secciones \ref{sec:Michelangelo} y \ref{sec:CUPID}.
    \item Bingbing Huang: Secciones \ref{sec:Michelangelo} y \ref{sec:CUPID}.
    \item Kerui Min: Traducción al chino.
    \item Prof. Qing Qu: Sección \ref{sec:diffusion_memorization}.
    \item Shengbang Peter Tong: Secciones \ref{sec:continuous} y \ref{sec:RAE-implement}.
    \item Chengyu Wang: Secciones \ref{sec:Michelangelo} y \ref{sec:CUPID}.
    \item Ziyang Robin Wu: Secciones \ref{sec:CLIP} y desarrollo del sitio web.
    \item Chun-Hsiao Daniel Yeh: Secciones \ref{sec:auto-encode-img-gen} y \ref{sec:cond_gen}.
    \item Brent Yi: Sección \ref{sec:egoallo}.
    \item Jingfeng Yang: Sección \ref{sec:CLIP}.
    \item Dr. Yaodong Yu: El capítulo \ref{ch:deep-networks} se basa en su tesis doctoral.
    \item Dr. Zibo Zhao: Sección \ref{sec:Michelangelo}.
\end{itemize}
También queremos agradecer al Dr. Kevin Murphy y al Dr. Bill Mark por sus extensos comentarios técnicos 
sobre el manuscrito; a Jan Cavel por haber contribuido una traducción no oficial al rumano, así como a 
Stephen Butterfill y Jeroen Van Goey por haber contribuido correcciones y modificaciones al manuscrito.

\end{document}
